\documentclass[11pt, a4paper]{article}
% ── Shared preamble for all RTOS lab documents ────────────────────────────────
% Usage: % ── Shared preamble for all RTOS lab documents ────────────────────────────────
% Usage: % ── Shared preamble for all RTOS lab documents ────────────────────────────────
% Usage: \input{../preamble}  (from each labNN/ directory)
% ──────────────────────────────────────────────────────────────────────────────

% ── Packages ──────────────────────────────────────────────────────────────────
\usepackage[a4paper, top=2.5cm, bottom=2.5cm, left=2.5cm, right=2.5cm, headheight=28pt]{geometry}
\usepackage[utf8]{inputenc}
\usepackage[T1]{fontenc}
\usepackage{lmodern}
\usepackage{booktabs}
\usepackage{array}
\usepackage{tabularx}
\usepackage{longtable}
\usepackage{multirow}
\usepackage{colortbl}
\usepackage{xcolor}
\usepackage{titlesec}
\usepackage{enumitem}
\usepackage{hyperref}
\usepackage{fancyhdr}
\usepackage{parskip}
\usepackage{listings}
\usepackage[most]{tcolorbox}
\usepackage{fontawesome5}
\usepackage{microtype}
\usepackage{graphicx}
\usepackage{amsmath}

% ── Colors (KMUTNB brand) ──────────────────────────────────────────────────────
\definecolor{kmutnbBlue}{RGB}{0,56,116}
\definecolor{kmutnbGold}{RGB}{186,146,36}
\definecolor{lightGray}{RGB}{245,245,245}
\definecolor{midGray}{RGB}{200,200,200}
\definecolor{codeGray}{RGB}{248,248,248}
\definecolor{codeFrame}{RGB}{220,220,220}
\definecolor{tipteal}{RGB}{0,105,92}
\definecolor{warnAmber}{RGB}{121,85,0}
\definecolor{checkGreen}{RGB}{27,94,32}

% ── Section formatting ─────────────────────────────────────────────────────────
\titleformat{\section}
  {\large\bfseries\color{kmutnbBlue}}
  {\thesection.}{0.5em}{}
  [\color{kmutnbGold}\titlerule]

\titleformat{\subsection}
  {\normalsize\bfseries\color{kmutnbBlue}}
  {\thesubsection.}{0.5em}{}

\titleformat{\subsubsection}
  {\small\bfseries\color{kmutnbBlue}}
  {\thesubsubsection.}{0.5em}{}

\titlespacing*{\section}{0pt}{1.4ex plus .2ex}{0.6ex plus .1ex}
\titlespacing*{\subsection}{0pt}{1.0ex plus .2ex}{0.4ex plus .1ex}

% ── Listings (C and bash) ──────────────────────────────────────────────────────
\lstdefinestyle{cstyle}{
  language        = C,
  basicstyle      = \ttfamily\small,
  keywordstyle    = \bfseries\color{kmutnbBlue},
  commentstyle    = \itshape\color{gray},
  stringstyle     = \color{tipteal},
  numberstyle     = \tiny\color{midGray},
  numbers         = left,
  numbersep       = 6pt,
  stepnumber      = 1,
  backgroundcolor = \color{codeGray},
  frame           = single,
  framerule       = 0.4pt,
  rulecolor       = \color{codeFrame},
  breaklines      = true,
  showstringspaces= false,
  tabsize         = 4,
  xleftmargin     = 1.5em,
  xrightmargin    = 0.5em,
  aboveskip       = 0.8em,
  belowskip       = 0.8em,
}

\lstdefinestyle{bashstyle}{
  language        = bash,
  basicstyle      = \ttfamily\small,
  keywordstyle    = \color{kmutnbBlue},
  commentstyle    = \itshape\color{gray},
  backgroundcolor = \color{black!5},
  frame           = single,
  framerule       = 0.4pt,
  rulecolor       = \color{codeFrame},
  breaklines      = true,
  showstringspaces= false,
  xleftmargin     = 1.5em,
  xrightmargin    = 0.5em,
  aboveskip       = 0.6em,
  belowskip       = 0.6em,
}

\lstdefinestyle{textstyle}{
  basicstyle      = \ttfamily\footnotesize,
  backgroundcolor = \color{black!4},
  frame           = single,
  framerule       = 0.4pt,
  rulecolor       = \color{codeFrame},
  breaklines      = true,
  xleftmargin     = 1.5em,
  xrightmargin    = 0.5em,
  aboveskip       = 0.6em,
  belowskip       = 0.6em,
}

% ── Callout boxes ──────────────────────────────────────────────────────────────
\tcbuselibrary{skins, breakable}

\newtcolorbox{tipbox}[1][]{
  enhanced, breakable,
  colback=tipteal!8, colframe=tipteal, coltitle=white,
  fonttitle=\bfseries\small, title={\faLightbulb\quad Tip#1},
  left=6pt, right=6pt, top=4pt, bottom=4pt,
  boxrule=0.8pt, arc=3pt,
}

\newtcolorbox{warnbox}[1][]{
  enhanced, breakable,
  colback=warnAmber!8, colframe=kmutnbGold, coltitle=white,
  fonttitle=\bfseries\small, title={\faExclamationTriangle\quad Note#1},
  left=6pt, right=6pt, top=4pt, bottom=4pt,
  boxrule=0.8pt, arc=3pt,
}

\newtcolorbox{checkpointbox}[1][]{
  enhanced,
  colback=kmutnbBlue!6, colframe=kmutnbBlue, coltitle=white,
  fonttitle=\bfseries\small, title={\faCheckCircle\quad Checkpoint#1},
  left=6pt, right=6pt, top=4pt, bottom=4pt,
  boxrule=0.8pt, arc=3pt,
}

\newtcolorbox{questionbox}[1][]{
  enhanced, breakable,
  colback=kmutnbGold!10, colframe=kmutnbGold, coltitle=white,
  fonttitle=\bfseries\small, title={\faQuestion\quad Question#1},
  left=6pt, right=6pt, top=4pt, bottom=4pt,
  boxrule=0.8pt, arc=3pt,
}

% ── Checkbox list ──────────────────────────────────────────────────────────────
\newlist{checklist}{itemize}{1}
\setlist[checklist]{label=\faSquare[regular], leftmargin=1.5em, noitemsep}

% ── Misc helpers ───────────────────────────────────────────────────────────────
\newcommand{\file}[1]{\texttt{\small #1}}
\newcommand{\api}[1]{\texttt{\small #1}}

% ── Title block macro ─────────────────────────────────────────────────────────
% Usage: \labtitle{03}{Aperiodic Servers \& Producer--Consumer}{2}{CLO 1 \& CLO 2}
\newcommand{\labtitle}[4]{%
  \begin{center}
    {\color{kmutnbBlue}\rule{\linewidth}{2pt}}\\[6pt]
    {\LARGE\bfseries\color{kmutnbBlue} Laboratory #1}\\[4pt]
    {\large\bfseries #2}\\[2pt]
    {\normalsize Real-Time Operating Systems \enspace\textbullet\enspace
     M.Eng.\ ECE \enspace\textbullet\enspace KMUTNB}\\[8pt]
    {\color{kmutnbGold}\rule{\linewidth}{1pt}}
  \end{center}
  \vspace{0.3cm}
  \begin{tabular}{@{}p{3.8cm} p{10cm}@{}}
    \textbf{Platform}       & NXP FRDM-MCXN236 (ARM Cortex-M33 @ 150\,MHz) \\[2pt]
    \textbf{RTOS}           & FreeRTOS v10.6 \\[2pt]
    \textbf{Toolchain}      & ARM GNU Toolchain (\texttt{arm-none-eabi-gcc}) + Make \\[2pt]
    \textbf{Estimated time} & #3 hours \\[2pt]
    \textbf{CLO mapping}    & #4 \\
  \end{tabular}
  \vspace{0.4cm}
}
  (from each labNN/ directory)
% ──────────────────────────────────────────────────────────────────────────────

% ── Packages ──────────────────────────────────────────────────────────────────
\usepackage[a4paper, top=2.5cm, bottom=2.5cm, left=2.5cm, right=2.5cm, headheight=28pt]{geometry}
\usepackage[utf8]{inputenc}
\usepackage[T1]{fontenc}
\usepackage{lmodern}
\usepackage{booktabs}
\usepackage{array}
\usepackage{tabularx}
\usepackage{longtable}
\usepackage{multirow}
\usepackage{colortbl}
\usepackage{xcolor}
\usepackage{titlesec}
\usepackage{enumitem}
\usepackage{hyperref}
\usepackage{fancyhdr}
\usepackage{parskip}
\usepackage{listings}
\usepackage[most]{tcolorbox}
\usepackage{fontawesome5}
\usepackage{microtype}
\usepackage{graphicx}
\usepackage{amsmath}

% ── Colors (KMUTNB brand) ──────────────────────────────────────────────────────
\definecolor{kmutnbBlue}{RGB}{0,56,116}
\definecolor{kmutnbGold}{RGB}{186,146,36}
\definecolor{lightGray}{RGB}{245,245,245}
\definecolor{midGray}{RGB}{200,200,200}
\definecolor{codeGray}{RGB}{248,248,248}
\definecolor{codeFrame}{RGB}{220,220,220}
\definecolor{tipteal}{RGB}{0,105,92}
\definecolor{warnAmber}{RGB}{121,85,0}
\definecolor{checkGreen}{RGB}{27,94,32}

% ── Section formatting ─────────────────────────────────────────────────────────
\titleformat{\section}
  {\large\bfseries\color{kmutnbBlue}}
  {\thesection.}{0.5em}{}
  [\color{kmutnbGold}\titlerule]

\titleformat{\subsection}
  {\normalsize\bfseries\color{kmutnbBlue}}
  {\thesubsection.}{0.5em}{}

\titleformat{\subsubsection}
  {\small\bfseries\color{kmutnbBlue}}
  {\thesubsubsection.}{0.5em}{}

\titlespacing*{\section}{0pt}{1.4ex plus .2ex}{0.6ex plus .1ex}
\titlespacing*{\subsection}{0pt}{1.0ex plus .2ex}{0.4ex plus .1ex}

% ── Listings (C and bash) ──────────────────────────────────────────────────────
\lstdefinestyle{cstyle}{
  language        = C,
  basicstyle      = \ttfamily\small,
  keywordstyle    = \bfseries\color{kmutnbBlue},
  commentstyle    = \itshape\color{gray},
  stringstyle     = \color{tipteal},
  numberstyle     = \tiny\color{midGray},
  numbers         = left,
  numbersep       = 6pt,
  stepnumber      = 1,
  backgroundcolor = \color{codeGray},
  frame           = single,
  framerule       = 0.4pt,
  rulecolor       = \color{codeFrame},
  breaklines      = true,
  showstringspaces= false,
  tabsize         = 4,
  xleftmargin     = 1.5em,
  xrightmargin    = 0.5em,
  aboveskip       = 0.8em,
  belowskip       = 0.8em,
}

\lstdefinestyle{bashstyle}{
  language        = bash,
  basicstyle      = \ttfamily\small,
  keywordstyle    = \color{kmutnbBlue},
  commentstyle    = \itshape\color{gray},
  backgroundcolor = \color{black!5},
  frame           = single,
  framerule       = 0.4pt,
  rulecolor       = \color{codeFrame},
  breaklines      = true,
  showstringspaces= false,
  xleftmargin     = 1.5em,
  xrightmargin    = 0.5em,
  aboveskip       = 0.6em,
  belowskip       = 0.6em,
}

\lstdefinestyle{textstyle}{
  basicstyle      = \ttfamily\footnotesize,
  backgroundcolor = \color{black!4},
  frame           = single,
  framerule       = 0.4pt,
  rulecolor       = \color{codeFrame},
  breaklines      = true,
  xleftmargin     = 1.5em,
  xrightmargin    = 0.5em,
  aboveskip       = 0.6em,
  belowskip       = 0.6em,
}

% ── Callout boxes ──────────────────────────────────────────────────────────────
\tcbuselibrary{skins, breakable}

\newtcolorbox{tipbox}[1][]{
  enhanced, breakable,
  colback=tipteal!8, colframe=tipteal, coltitle=white,
  fonttitle=\bfseries\small, title={\faLightbulb\quad Tip#1},
  left=6pt, right=6pt, top=4pt, bottom=4pt,
  boxrule=0.8pt, arc=3pt,
}

\newtcolorbox{warnbox}[1][]{
  enhanced, breakable,
  colback=warnAmber!8, colframe=kmutnbGold, coltitle=white,
  fonttitle=\bfseries\small, title={\faExclamationTriangle\quad Note#1},
  left=6pt, right=6pt, top=4pt, bottom=4pt,
  boxrule=0.8pt, arc=3pt,
}

\newtcolorbox{checkpointbox}[1][]{
  enhanced,
  colback=kmutnbBlue!6, colframe=kmutnbBlue, coltitle=white,
  fonttitle=\bfseries\small, title={\faCheckCircle\quad Checkpoint#1},
  left=6pt, right=6pt, top=4pt, bottom=4pt,
  boxrule=0.8pt, arc=3pt,
}

\newtcolorbox{questionbox}[1][]{
  enhanced, breakable,
  colback=kmutnbGold!10, colframe=kmutnbGold, coltitle=white,
  fonttitle=\bfseries\small, title={\faQuestion\quad Question#1},
  left=6pt, right=6pt, top=4pt, bottom=4pt,
  boxrule=0.8pt, arc=3pt,
}

% ── Checkbox list ──────────────────────────────────────────────────────────────
\newlist{checklist}{itemize}{1}
\setlist[checklist]{label=\faSquare[regular], leftmargin=1.5em, noitemsep}

% ── Misc helpers ───────────────────────────────────────────────────────────────
\newcommand{\file}[1]{\texttt{\small #1}}
\newcommand{\api}[1]{\texttt{\small #1}}

% ── Title block macro ─────────────────────────────────────────────────────────
% Usage: \labtitle{03}{Aperiodic Servers \& Producer--Consumer}{2}{CLO 1 \& CLO 2}
\newcommand{\labtitle}[4]{%
  \begin{center}
    {\color{kmutnbBlue}\rule{\linewidth}{2pt}}\\[6pt]
    {\LARGE\bfseries\color{kmutnbBlue} Laboratory #1}\\[4pt]
    {\large\bfseries #2}\\[2pt]
    {\normalsize Real-Time Operating Systems \enspace\textbullet\enspace
     M.Eng.\ ECE \enspace\textbullet\enspace KMUTNB}\\[8pt]
    {\color{kmutnbGold}\rule{\linewidth}{1pt}}
  \end{center}
  \vspace{0.3cm}
  \begin{tabular}{@{}p{3.8cm} p{10cm}@{}}
    \textbf{Platform}       & NXP FRDM-MCXN236 (ARM Cortex-M33 @ 150\,MHz) \\[2pt]
    \textbf{RTOS}           & FreeRTOS v10.6 \\[2pt]
    \textbf{Toolchain}      & ARM GNU Toolchain (\texttt{arm-none-eabi-gcc}) + Make \\[2pt]
    \textbf{Estimated time} & #3 hours \\[2pt]
    \textbf{CLO mapping}    & #4 \\
  \end{tabular}
  \vspace{0.4cm}
}
  (from each labNN/ directory)
% ──────────────────────────────────────────────────────────────────────────────

% ── Packages ──────────────────────────────────────────────────────────────────
\usepackage[a4paper, top=2.5cm, bottom=2.5cm, left=2.5cm, right=2.5cm, headheight=28pt]{geometry}
\usepackage[utf8]{inputenc}
\usepackage[T1]{fontenc}
\usepackage{lmodern}
\usepackage{booktabs}
\usepackage{array}
\usepackage{tabularx}
\usepackage{longtable}
\usepackage{multirow}
\usepackage{colortbl}
\usepackage{xcolor}
\usepackage{titlesec}
\usepackage{enumitem}
\usepackage{hyperref}
\usepackage{fancyhdr}
\usepackage{parskip}
\usepackage{listings}
\usepackage[most]{tcolorbox}
\usepackage{fontawesome5}
\usepackage{microtype}
\usepackage{graphicx}
\usepackage{amsmath}

% ── Colors (KMUTNB brand) ──────────────────────────────────────────────────────
\definecolor{kmutnbBlue}{RGB}{0,56,116}
\definecolor{kmutnbGold}{RGB}{186,146,36}
\definecolor{lightGray}{RGB}{245,245,245}
\definecolor{midGray}{RGB}{200,200,200}
\definecolor{codeGray}{RGB}{248,248,248}
\definecolor{codeFrame}{RGB}{220,220,220}
\definecolor{tipteal}{RGB}{0,105,92}
\definecolor{warnAmber}{RGB}{121,85,0}
\definecolor{checkGreen}{RGB}{27,94,32}

% ── Section formatting ─────────────────────────────────────────────────────────
\titleformat{\section}
  {\large\bfseries\color{kmutnbBlue}}
  {\thesection.}{0.5em}{}
  [\color{kmutnbGold}\titlerule]

\titleformat{\subsection}
  {\normalsize\bfseries\color{kmutnbBlue}}
  {\thesubsection.}{0.5em}{}

\titleformat{\subsubsection}
  {\small\bfseries\color{kmutnbBlue}}
  {\thesubsubsection.}{0.5em}{}

\titlespacing*{\section}{0pt}{1.4ex plus .2ex}{0.6ex plus .1ex}
\titlespacing*{\subsection}{0pt}{1.0ex plus .2ex}{0.4ex plus .1ex}

% ── Listings (C and bash) ──────────────────────────────────────────────────────
\lstdefinestyle{cstyle}{
  language        = C,
  basicstyle      = \ttfamily\small,
  keywordstyle    = \bfseries\color{kmutnbBlue},
  commentstyle    = \itshape\color{gray},
  stringstyle     = \color{tipteal},
  numberstyle     = \tiny\color{midGray},
  numbers         = left,
  numbersep       = 6pt,
  stepnumber      = 1,
  backgroundcolor = \color{codeGray},
  frame           = single,
  framerule       = 0.4pt,
  rulecolor       = \color{codeFrame},
  breaklines      = true,
  showstringspaces= false,
  tabsize         = 4,
  xleftmargin     = 1.5em,
  xrightmargin    = 0.5em,
  aboveskip       = 0.8em,
  belowskip       = 0.8em,
}

\lstdefinestyle{bashstyle}{
  language        = bash,
  basicstyle      = \ttfamily\small,
  keywordstyle    = \color{kmutnbBlue},
  commentstyle    = \itshape\color{gray},
  backgroundcolor = \color{black!5},
  frame           = single,
  framerule       = 0.4pt,
  rulecolor       = \color{codeFrame},
  breaklines      = true,
  showstringspaces= false,
  xleftmargin     = 1.5em,
  xrightmargin    = 0.5em,
  aboveskip       = 0.6em,
  belowskip       = 0.6em,
}

\lstdefinestyle{textstyle}{
  basicstyle      = \ttfamily\footnotesize,
  backgroundcolor = \color{black!4},
  frame           = single,
  framerule       = 0.4pt,
  rulecolor       = \color{codeFrame},
  breaklines      = true,
  xleftmargin     = 1.5em,
  xrightmargin    = 0.5em,
  aboveskip       = 0.6em,
  belowskip       = 0.6em,
}

% ── Callout boxes ──────────────────────────────────────────────────────────────
\tcbuselibrary{skins, breakable}

\newtcolorbox{tipbox}[1][]{
  enhanced, breakable,
  colback=tipteal!8, colframe=tipteal, coltitle=white,
  fonttitle=\bfseries\small, title={\faLightbulb\quad Tip#1},
  left=6pt, right=6pt, top=4pt, bottom=4pt,
  boxrule=0.8pt, arc=3pt,
}

\newtcolorbox{warnbox}[1][]{
  enhanced, breakable,
  colback=warnAmber!8, colframe=kmutnbGold, coltitle=white,
  fonttitle=\bfseries\small, title={\faExclamationTriangle\quad Note#1},
  left=6pt, right=6pt, top=4pt, bottom=4pt,
  boxrule=0.8pt, arc=3pt,
}

\newtcolorbox{checkpointbox}[1][]{
  enhanced,
  colback=kmutnbBlue!6, colframe=kmutnbBlue, coltitle=white,
  fonttitle=\bfseries\small, title={\faCheckCircle\quad Checkpoint#1},
  left=6pt, right=6pt, top=4pt, bottom=4pt,
  boxrule=0.8pt, arc=3pt,
}

\newtcolorbox{questionbox}[1][]{
  enhanced, breakable,
  colback=kmutnbGold!10, colframe=kmutnbGold, coltitle=white,
  fonttitle=\bfseries\small, title={\faQuestion\quad Question#1},
  left=6pt, right=6pt, top=4pt, bottom=4pt,
  boxrule=0.8pt, arc=3pt,
}

% ── Checkbox list ──────────────────────────────────────────────────────────────
\newlist{checklist}{itemize}{1}
\setlist[checklist]{label=\faSquare[regular], leftmargin=1.5em, noitemsep}

% ── Misc helpers ───────────────────────────────────────────────────────────────
\newcommand{\file}[1]{\texttt{\small #1}}
\newcommand{\api}[1]{\texttt{\small #1}}

% ── Title block macro ─────────────────────────────────────────────────────────
% Usage: \labtitle{03}{Aperiodic Servers \& Producer--Consumer}{2}{CLO 1 \& CLO 2}
\newcommand{\labtitle}[4]{%
  \begin{center}
    {\color{kmutnbBlue}\rule{\linewidth}{2pt}}\\[6pt]
    {\LARGE\bfseries\color{kmutnbBlue} Laboratory #1}\\[4pt]
    {\large\bfseries #2}\\[2pt]
    {\normalsize Real-Time Operating Systems \enspace\textbullet\enspace
     M.Eng.\ ECE \enspace\textbullet\enspace KMUTNB}\\[8pt]
    {\color{kmutnbGold}\rule{\linewidth}{1pt}}
  \end{center}
  \vspace{0.3cm}
  \begin{tabular}{@{}p{3.8cm} p{10cm}@{}}
    \textbf{Platform}       & NXP FRDM-MCXN236 (ARM Cortex-M33 @ 150\,MHz) \\[2pt]
    \textbf{RTOS}           & FreeRTOS v10.6 \\[2pt]
    \textbf{Toolchain}      & ARM GNU Toolchain (\texttt{arm-none-eabi-gcc}) + Make \\[2pt]
    \textbf{Estimated time} & #3 hours \\[2pt]
    \textbf{CLO mapping}    & #4 \\
  \end{tabular}
  \vspace{0.4cm}
}


% ── Header / Footer ────────────────────────────────────────────────────────────
\pagestyle{fancy}
\fancyhf{}
\fancyhead[L]{\small\color{kmutnbBlue}\textbf{KMUTNB -- Faculty of Engineering}}
\fancyhead[R]{\small\color{kmutnbBlue}Dept.\ of Electrical and Computer Engineering}
\fancyfoot[L]{\small\color{kmutnbBlue}Real-Time Operating Systems -- M.Eng.\ ECE}
\fancyfoot[C]{\small\thepage}
\fancyfoot[R]{\small\color{kmutnbBlue}Lab 06}
\renewcommand{\headrulewidth}{0.4pt}
\renewcommand{\headrule}{\hbox to\headwidth{%
  \color{kmutnbGold}\leaders\hrule height\headrulewidth\hfill}}
\renewcommand{\footrulewidth}{0.4pt}
\renewcommand{\footrule}{\hbox to\headwidth{%
  \color{midGray}\leaders\hrule height\footrulewidth\hfill}}

\hypersetup{
  colorlinks = true,
  linkcolor  = kmutnbBlue,
  urlcolor   = kmutnbBlue,
  citecolor  = kmutnbBlue,
  pdftitle   = {Lab 06 -- WCET Measurement with DWT \& Interrupt Latency},
  pdfauthor  = {KMUTNB RTOS Course},
}

% ══════════════════════════════════════════════════════════════════════════════
\begin{document}

\labtitle{06}{WCET Measurement with DWT \& Interrupt Latency}{3--4}{CLO 3 \& CLO 4}

% ── Objectives ────────────────────────────────────────────────────────────────
\section{Objectives}

By completing this lab you will be able to:

\begin{enumerate}[noitemsep, leftmargin=*]
  \item \textbf{Enable and read} the DWT cycle counter on the ARM Cortex-M33 at
        150\,MHz.
  \item \textbf{Measure} best-case, average, and worst-case execution time for a
        data-processing function over 1\,000 runs.
  \item \textbf{Apply} a WCET safety margin and use the result as $C_i$ in
        schedulability analysis.
  \item \textbf{Measure} interrupt latency from IRQ assertion to first ISR instruction
        using an oscilloscope or logic analyser.
  \item \textbf{Compare} measured interrupt latency with and without a FreeRTOS critical
        section active.
\end{enumerate}

% ── Prerequisites ─────────────────────────────────────────────────────────────
\section{Prerequisites}

\begin{checklist}
  \item Lab 03 complete.
  \item ARM GNU Toolchain and \texttt{pyocd} working.
  \item Oscilloscope or logic analyser with $\geq 4$ channels (for Part~B). A probe on
        GPIO pins is sufficient.
  \item FRDM-MCXN236 board --- no other hardware required for Part~A.
\end{checklist}

% ── Background ────────────────────────────────────────────────────────────────
\section{Background}

\subsection{Why Measure WCET?}

The Response-Time Analysis from Week 3 requires a bound $C_i$ on task execution time. In
practice, $C_i$ is obtained by:

\begin{enumerate}[noitemsep, leftmargin=*]
  \item \textbf{Static analysis} --- examine all paths, compute worst-case bound.
        Requires special tools (aiT, Bound-T).
  \item \textbf{Measurement-based} --- run the function many times with controlled
        inputs, take the maximum plus a safety margin.
\end{enumerate}

This lab uses measurement-based WCET --- practical for most systems below SIL-3/ASIL-D.

\subsection{DWT Cycle Counter}

The ARM Data Watchpoint and Trace (DWT) unit has a 32-bit free-running cycle counter.
Steps to enable:

\begin{lstlisting}[style=cstyle]
CoreDebug->DEMCR |= CoreDebug_DEMCR_TRCENA_Msk;  /* enable DWT */
DWT->CYCCNT = 0;
DWT->CTRL   |= DWT_CTRL_CYCCNTENA_Msk;            /* start counting */
\end{lstlisting}

Read: \texttt{uint32\_t cycles = DWT->CYCCNT;}

At 150\,MHz, 1 cycle = 6.67\,ns. The counter rolls over after $2^{32}$ cycles
$\approx 28.6$\,seconds.

% ── Project Structure ─────────────────────────────────────────────────────────
\section{Project Structure}

\begin{lstlisting}[style=textstyle]
lab06_wcet_dwt/
  Makefile
  linker_MCXN236.ld
  src/
    main.c              <- DWT init, task creation
    dwt.h               <- DWT enable/read macros
    workload.c/.h       <- the functions to be measured
    latency_test.c/.h   <- Part B: GPIO ISR latency test
    FreeRTOSConfig.h
    uart.h / uart.c
  FreeRTOS-Kernel/
\end{lstlisting}

% ── Part A ─────────────────────────────────────────────────────────────────────
\section{Part A --- WCET Measurement}

\subsection{Step 1 --- Enable DWT}

In \file{dwt.h}:

\begin{lstlisting}[style=cstyle]
#ifndef DWT_H
#define DWT_H
#include "core_cm33.h"

static inline void DWT_Enable(void)
{
    CoreDebug->DEMCR |= CoreDebug_DEMCR_TRCENA_Msk;
    DWT->CYCCNT = 0;
    DWT->CTRL  |= DWT_CTRL_CYCCNTENA_Msk;
    __DSB();
    __ISB();
}

static inline uint32_t DWT_GetCycles(void)
{
    return DWT->CYCCNT;
}

#define DWT_START()  uint32_t _dwt_start = DWT_GetCycles()
#define DWT_STOP()   uint32_t _dwt_elapsed = DWT_GetCycles() - _dwt_start
#endif
\end{lstlisting}

\subsection{Step 2 --- Implement the Workload Functions}

Three functions of increasing complexity in \file{workload.c}:

\begin{lstlisting}[style=cstyle]
/* Function 1: bubble sort -- data-dependent WCET */
void bubble_sort(int16_t *arr, uint16_t n)
{
    for (uint16_t i = 0; i < n - 1; i++)
        for (uint16_t j = 0; j < n - i - 1; j++)
            if (arr[j] > arr[j+1]) {
                int16_t tmp = arr[j];
                arr[j]   = arr[j+1];
                arr[j+1] = tmp;
            }
}

/* Function 2: 16-tap FIR filter -- fixed WCET */
int32_t fir_filter(const int16_t *x, const int16_t *h, uint16_t n)
{
    int32_t acc = 0;
    for (uint16_t i = 0; i < n; i++)
        acc += (int32_t)x[i] * h[i];
    return acc >> 15;
}

/* Function 3: 32-bit CRC computation -- data-dependent */
uint32_t crc32(const uint8_t *data, uint32_t len)
{
    uint32_t crc = 0xFFFFFFFFU;
    while (len--) {
        crc ^= *data++;
        for (int k = 0; k < 8; k++)
            crc = (crc >> 1) ^ (0xEDB88320U & -(crc & 1));
    }
    return ~crc;
}
\end{lstlisting}

\subsection{Step 3 --- Measurement Task}

\begin{lstlisting}[style=cstyle]
#define NUM_RUNS     1000
#define ARRAY_SIZE   100

void vWcetMeasureTask(void *pv)
{
    static int16_t sorted[ARRAY_SIZE], reverse[ARRAY_SIZE], arr[ARRAY_SIZE];
    static int16_t fir_h[16] = {1,2,3,4,5,6,7,8,8,7,6,5,4,3,2,1};
    uint32_t best, worst, sum, cycles;

    for (int i = 0; i < ARRAY_SIZE; i++) {
        sorted[i]  = (int16_t)i;
        reverse[i] = (int16_t)(ARRAY_SIZE - 1 - i);
    }

    uart_printf("\r\n=== bubble_sort(%d elements) ===\r\n", ARRAY_SIZE);

    /* Worst case: reverse sorted, NUM_RUNS times */
    best = UINT32_MAX; worst = 0; sum = 0;
    for (int r = 0; r < NUM_RUNS; r++) {
        memcpy(arr, reverse, sizeof(reverse));
        DWT_START();
        bubble_sort(arr, ARRAY_SIZE);
        DWT_STOP();
        cycles = _dwt_elapsed;
        if (cycles < best) best = cycles;
        if (cycles > worst) worst = cycles;
        sum += cycles;
    }
    uart_printf("  best=%lu  worst=%lu  avg=%lu cycles\r\n",
                best, worst, sum / NUM_RUNS);
    uart_printf("  WCET with 20%% margin: %lu cycles  (%.2f us)\r\n",
                (uint32_t)(worst * 1.2f), worst * 1.2f / 150.0f);

    uart_printf("\r\nMeasurement complete.\r\n");
    vTaskDelete(NULL);
}
\end{lstlisting}

\subsection{Step 4 --- Build, Flash, and Record Results}

\begin{lstlisting}[style=bashstyle]
make setup && make && make flash
\end{lstlisting}

\begin{checkpointbox}[~A]
  Terminal output with all three function measurements.
\end{checkpointbox}

\begin{questionbox}[~A1]
  Bubble sort has the same code path for all reverse-sorted inputs. Why does
  \texttt{worst} vary across the 1\,000 runs even with identical input? Name two
  hardware effects that cause this variability.
\end{questionbox}

\begin{questionbox}[~A2]
  The FIR filter is data-independent (no branches). Does it show variability across
  runs? What is the ratio of worst to best cycles? Explain.
\end{questionbox}

\begin{questionbox}[~A3]
  You will use \texttt{bubble\_sort} as a task in Lab~07 with period $T = 50$\,ms.
  Using the WCET + 20\% margin as $C_i$: compute the utilisation $C_i / T_i$. At what
  maximum number of identical \texttt{bubble\_sort} tasks does the total utilisation
  exceed the Liu \& Layland bound for RMS?
\end{questionbox}

% ── Part B ─────────────────────────────────────────────────────────────────────
\section{Part B --- Interrupt Latency Measurement}

\subsection{Step 1 --- Set Up the GPIO Toggling ISR}

Wire a jumper from \texttt{P1\_0} (output) to \texttt{P1\_1} (input, interrupt on rising
edge). The ISR immediately toggles \texttt{P1\_2} (output to oscilloscope probe).

\begin{lstlisting}[style=cstyle]
/* latency_test.c */
void GPIO1_IRQHandler(void)
{
    GPIO_ClearPinsInterruptFlags(GPIO1, 1u << 1);
    GPIO_TogglePinsOutput(GPIO1, 1u << 2);   /* respond ASAP */
}
\end{lstlisting}

\subsection{Step 2 --- Configure GPIO and NVIC}

\begin{lstlisting}[style=cstyle]
/* Input: P1_1 rising edge interrupt */
gpio_pin_config_t in_cfg  = { kGPIO_DigitalInput,  0, kGPIO_IntRisingEdge };
GPIO_PinInit(GPIO1, 1, &in_cfg);
GPIO_EnableInterrupts(GPIO1, 1u << 1);
EnableIRQ(GPIO1_IRQn);

/* Output: P1_2 */
gpio_pin_config_t out_cfg = { kGPIO_DigitalOutput, 0, kGPIO_NoIntmode };
GPIO_PinInit(GPIO1, 2, &out_cfg);
\end{lstlisting}

\subsection{Step 3 --- Drive the Input and Measure}

Three measurement options:

\begin{itemize}[noitemsep, leftmargin=*]
  \item \textbf{Option A --- Logic analyser:} toggle \texttt{P1\_0} at 1\,kHz from a
        software task; measure \texttt{P1\_1}-rising to \texttt{P1\_2}-rising.
  \item \textbf{Option B --- Oscilloscope:} probe \texttt{P1\_1} (CH1) and
        \texttt{P1\_2} (CH2); trigger on CH1 rising; measure propagation delay.
  \item \textbf{Option C --- DWT inside the ISR:} record \texttt{DWT->CYCCNT} as the
        first ISR instruction and compare against the cycle count set just before
        toggling \texttt{P1\_0}.
\end{itemize}

\subsection{Step 4 --- Measure With and Without Critical Section}

Repeat the measurement with a FreeRTOS critical section active in a background task:

\begin{lstlisting}[style=cstyle]
void vCriticalSectionTask(void *pv)
{
    for (;;) {
        taskENTER_CRITICAL();
        volatile uint32_t i = 0;
        while (i++ < 150000) {}   /* ~1 ms at 150 MHz */
        taskEXIT_CRITICAL();
        vTaskDelay(pdMS_TO_TICKS(5));
    }
}
\end{lstlisting}

\begin{checkpointbox}[~B]
  Table or oscilloscope screenshot showing latency with and without the critical
  section.
\end{checkpointbox}

\begin{questionbox}[~B1]
  Your measured best-case latency should be 12--18 cycles. Convert this to nanoseconds
  at 150\,MHz. How does this compare to the theoretical minimum from the ARM
  Cortex-M33 TRM (exception entry stacking)?
\end{questionbox}

\begin{questionbox}[~B2]
  With the critical section active, what is the worst-case interrupt latency? Where does
  the extra latency come from? What FreeRTOS configuration parameter controls this?
\end{questionbox}

\begin{questionbox}[~B3]
  The ISR priority is set to
  \texttt{configLIBRARY\_MAX\_SYSCALL\_INTERRUPT\_PRIORITY - 1}. Why can this ISR still
  be deferred by \texttt{taskENTER\_CRITICAL}? What priority would you set to make the
  ISR \textbf{non-maskable} by FreeRTOS?
\end{questionbox}

% ── Part C ─────────────────────────────────────────────────────────────────────
\section{Part C --- Full Task WCET in Context}

\subsection{Step 1 --- Measure Task Response Time End-to-End}

Create a periodic task that runs \texttt{bubble\_sort} on random data every 50\,ms.
Use DWT to measure from task entry to task return, with two background tasks also active
(priorities 1 and 2):

\begin{lstlisting}[style=cstyle]
void vSortTask(void *pv)
{
    static int16_t arr[ARRAY_SIZE];
    TickType_t xLastWake = xTaskGetTickCount();
    for (;;) {
        vTaskDelayUntil(&xLastWake, pdMS_TO_TICKS(50));
        for (int i = 0; i < ARRAY_SIZE; i++) arr[i] = (int16_t)(rand() % 1000);
        DWT_START();
        bubble_sort(arr, ARRAY_SIZE);
        DWT_STOP();
        uart_printf("[SORT] %lu cycles  (%.2f us)\r\n",
                    _dwt_elapsed, _dwt_elapsed / 150.0f);
    }
}
\end{lstlisting}

Compare the measured cycles to Part~A's isolated measurement.

\begin{questionbox}[~C1]
  Does running alongside other tasks increase the measured execution time? What hardware
  effect causes this? What does this tell you about measurement-based WCET methodology?
\end{questionbox}

% ── Deliverables ───────────────────────────────────────────────────────────────
\section{Deliverables}

\renewcommand{\arraystretch}{1.3}
\begin{tabular}{@{} c p{4.5cm} p{7cm} @{}}
  \toprule
  \textbf{\#} & \textbf{Item} & \textbf{Details} \\
  \midrule
  1 & Part A terminal output    & Full measurement table for all three functions \\
  2 & Part A analysis           & WCET + 20\% margin and schedulability (Question~A3) \\
  3 & Part B measurement        & Latency table or oscilloscope screenshot \\
  4 & Part B analysis           & With and without critical section (Question~B2) \\
  5 & Written answers           & Questions A1--A3, B1--B3, C1 \\
  6 & \file{dwt.h} header file  & Full source listing \\
  \bottomrule
\end{tabular}
\renewcommand{\arraystretch}{1}

% ── Key Concepts ───────────────────────────────────────────────────────────────
\section{Key Concepts Demonstrated}

\renewcommand{\arraystretch}{1.3}
\begin{tabular}{@{} p{7cm} p{6cm} @{}}
  \toprule
  \textbf{Concept} & \textbf{Where you saw it} \\
  \midrule
  DWT cycle counter enable and read         & Part A setup \\
  Best/average/worst-case measurement       & Part A \\
  WCET safety margin application            & Part A Question~A3 \\
  Interrupt latency: hardware stacking      & Part B \\
  BASEPRI masking and FreeRTOS critical section & Part B \\
  Context and cache effects on WCET         & Part C \\
  \bottomrule
\end{tabular}
\renewcommand{\arraystretch}{1}

% ── Reference ─────────────────────────────────────────────────────────────────
\section{Reference}

\renewcommand{\arraystretch}{1.3}
\begin{tabular}{@{} p{5.5cm} p{7.5cm} @{}}
  \toprule
  \textbf{Resource} & \textbf{Relevant sections} \\
  \midrule
  Wilhelm et al.\ (2008) &
    ``The worst-case execution time problem'' (\textit{ACM TECS}) \\
  ARM Cortex-M33 TRM & \S DWT, \S B2.3 exception entry \\
  NXP MCXN236 RM & \S CoreDebug, \S DWT \\
  FreeRTOS docs &
    \texttt{taskENTER\_CRITICAL}, \texttt{configMAX\_SYSCALL\_INTERRUPT\_PRIORITY} \\
  \bottomrule
\end{tabular}
\renewcommand{\arraystretch}{1}

\end{document}
